\section{Introduction}

The exponential growth of scientific literature has created a demand for text summarization methods to support scientists in prioritizing articles, extracting relevant information, and interpreting research findings. \ac{ATS} methods have evolved from statistical approaches to deep learning-based models, thus becoming increasingly sophisticated and reliable in capturing the essential parts of complex research articles. Although \ac{ATS} methods have been previously evaluated and described  \cite{zhang2025comprehensivesurveyprocessorientedautomatic,zhang2024systematicsurveytextsummarization}, only a few have focused on scientific literature \cite{ROHIL2022100058,xie2023surveybiomedicaltextsummarization}. Recent work has begun to benchmark \acp{LLM} on biomedical summarization, though only for a limited set of models and typically against abstract-based references \cite{chen2025benchmarking,celikten2025benchmarking,toprak2025benchmarking}. A systematic comparison from word-frequency methods to state-of-the-art \acp{LLM} is still lacking.

Over time, these methods have developed into distinct approaches that differ in how they identify or generate summary content. Early pre-neural summarization was mainly characterized by unsupervised extractive approaches that rank sentences by word or concept frequencies to identify relevant sentences. The field was initially shaped by Luhn's assumption that recurrent words signal importance   \cite{Luhn1958TheAC} and Edmundson's use of cue words, title words, and sentence position \cite{10.1145/321510.321519}. Later work adopting weighting \ac{TF-IDF} refined this idea by down-weighting frequent, low-specificity terms \cite{article}, laying the foundation for more advanced methods in scientific text summarization \cite{10.1145/1183614.1183701}. Graph-based methods judged the importance against the document's global structure rather than local word counts. TextRank, widely used in the biomedical domain, applies the PageRank algorithm to a sentence-similarity graph to select top-ranked sentences \cite{mihalcea-tarau-2004-textrank}. 

The advent of \ac{Seq2seq} frameworks shifted the field toward neural approaches capable of paraphrasing and condensing text using encoder-decoder architectures, originally implemented with \acp{RNN}, \ac{LSTM} networks, and \acp{GRU} \cite{afzal2020, almasoud2022}, while self-attention enabled parallel processing of sequences and modeling of long-range context  \cite{vaswani2023attentionneed}. This laid the foundation for transformer architectures that quickly gained popularity in performing a wide range of \ac{NLP} tasks. Building on \ac{BERT} \cite{devlin2019bertpretrainingdeepbidirectional}, several abstractive summarization models have emerged, including \ac{BART}, a denoising autoencoder for pretraining \ac{Seq2seq} models \cite{DBLP:journals/corr/abs-1910-13461} and its biomedical adaptations \cite{yuan-etal-2022-biobart, abinaya2024}; the unified \ac{T5} model \cite{JMLR:v21:20-074}; \ac{PEGASUS} \cite{zhang2020pegasuspretrainingextractedgapsentences}, including domain-specific PubMed variants; and Longformer-based models \cite{Beltagy2020Longformer} that can handle longer inputs \cite{steblianko2024}. These models are typically pretrained independently, while Longformer-based variants may additionally leverage \ac{RoBERTa}, an optimized version of \ac{BERT} trained on a larger corpus.   

More recently, the field of \ac{ATS} has quickly moved toward decoder-only architectures at the core of \acp{LLM}, which capture semantic relations with greater flexibility and specificity. \acp{LLM} can be classified as (i) general-purpose models, which apply broad domain knowledge across diverse \ac{NLP} tasks, such as the proprietary \ac{GPT} \cite{Radford2018ImprovingLU} and Claude \cite{bai2022constitutionalaiharmlessnessai} series and the open-source \ac{Llama} \cite{grattafiori2024llama3herdmodels} and Gemma \cite{gemmateam2025gemma3technicalreport} families; (ii) reasoning-oriented models, characterized by logical text understanding through iterative chain-of-thought processing and instruction tuning \cite{plaat2025multistepreasoninglargelanguage} such as DeepSeek-R1 \cite{wang2025reviewdeepseekmodelskey}, Qwen3 \cite{qwen3technicalreport}, and Mistral's Magistral \cite{mistralai2025magistral}, together with the reasoning tiers of the proprietary families (e.g., GPT-5 and Claude Opus 4); and (iii) domain-specific models, tailored for specialized tasks or specific scientific domains , such as the biomedical models OpenBioLLM-Llama3 \cite{OpenBioLLMs}, BioGPT \cite{10.1093/bib/bbac409}, and BioMistral \cite{labrak2024biomistral}, and SciLitLLM \cite{li2025scilitllmadaptllmsscientific} for scientific literature. The gold-standard dataset used for benchmarking and the complete set of models evaluated in this study are listed in Tables~\ref{tab:goldstandard} and \ref{tab:models} in the STAR Methods section.

Despite this range of approaches, these method families have not been evaluated together on biomedical literature, leaving little guidance on which to choose. This study addresses this gap by systematically evaluating 62 text summarization models, ranging from word-frequency methods to state-of-the-art \acp{LLM}, using a dataset of 1,000 biomedical abstracts and corresponding highlights sections as reference summaries. We provide recommendations for selecting appropriate tools to accelerate knowledge discovery in biomedical sciences and discuss the strengths and limitations of each approach.  